% Chapter 6

\chapter{CONCLUSION AND FUTURE WORK}\label{cha:conclusion}

\section{CONCLUSION}

This project set out to let the payload decide the route. Existing multi-objective UAV
planners build good Pareto fronts and then choose from them with a rule that knows
nothing about the cargo, or with weights an expert re-tunes for every mission. We
separated the two problems. Tier~1 builds an exact NAMOA* front over a
three-dimensional, wind-aware airspace lattice once per region and weather window, and
stores it with what is needed to tell when it has gone stale. Tier~2 turns four payload
tags into objective weights with an explicit formula, scores every route on the front
and returns the best one, or a HOLD when the weather or the battery rules the mission
out. The chosen route is exported as an ArduPilot mission and has been flown in SITL and
in a Gazebo city world.

The first results are encouraging, with clear caveats. The Tier~1 front was exact in
every benchmark scenario, while MOPSO found 38 of 115 routes on the largest. On that
front the four reference payloads received four different routes, three of them
non-extreme, in about 0.14\,ms each. Wind direction alone moved the energy of the
cheapest route by 27.5\% at a fixed speed, and in SITL the flown track stayed within
0.13\,m RMS of the plan. On the other hand, under a steady wind the front shrank and the
presets collapsed onto two routes, MOPSO ran at reduced settings, the ablation used
fixture data, and the risk weights and battery constants are still placeholders.

\section{FUTURE WORK}

The rest of Phase~I is a set of passes through the evaluation cycle. We will run the
full benchmark (four graphs, nine wind cases, MOPSO at 100 particles, 200 iterations and
30 repetitions) with per-repetition statistics and a significance test; fix the density
labels and the monotonicity grouping; study why the front collapses under steady wind,
testing other normalisations and weight multipliers; repeat the ablation and Tier~2
studies on live map data; compare Tier~2 with crowding distance, knee point, shortest
route and fixed weights over all 225 tag combinations; close the two dashboard defects;
and reduce the Tier~1 runtime on large graphs. On the simulation side, wind will be added
to the Gazebo world, the battery model will be calibrated per payload weight class, and
each preset will be flown under several wind cases with all flight metrics logged.

Phase~II moves the system onto a Pixhawk-equipped F450 quadcopter with a Raspberry
Pi~5 companion computer and a servo payload release. Tier~2 will run on board, and
outdoor flights will measure route adherence and the time from tag entry to mission
start. Further out, the two-tier idea could extend to several drones sharing one front,
to multi-drop routes and to in-flight replanning when the weather changes faster than
the refresh window.
